\documentclass[article]{jss}

\usepackage{amsmath}
\usepackage{amssymb}
\usepackage{amsfonts}
\usepackage{graphicx}
\usepackage{url}
\usepackage{hyperref}
\usepackage{booktabs}
\usepackage{color}
\usepackage{natbib}

%% Fallbacks so the paper also builds against a minimal jss.cls (the real JSS
%% class defines all of these). Harmless when the full class is present.
\providecommand{\proglang}[1]{\textsf{#1}}
\providecommand{\pkg}[1]{\textbf{#1}}
\providecommand{\code}[1]{\texttt{#1}}

%% JSS article information
\author{
  Peter Cotton\\
  Microprediction\\
  \email{peter.cotton@microprediction.com}\\
  \url{https://www.microprediction.org}
}

\title{\pkg{skaters}: Model First, Conform Last --- Composable Online Distributional Time-Series Forecasting in \proglang{Python} and \proglang{JavaScript}}

\Plaintitle{skaters: Model First, Conform Last}
\Shorttitle{\pkg{skaters}: Online Distributional Forecasting}

\newcommand{\skatersabstract}{
We present \pkg{skaters}, a zero-dependency library for online univariate
time-series forecasting in which \emph{every prediction is a full probability
distribution}, not a point. Models are built by composition: invertible
transforms chain together, ensembles nest, and a single distributional
\emph{leaf} sits at the bottom of every chain. The leaf is a pluggable
proper-scoring-rule optimiser, which lets us separate two concerns that the
forecasting and conformal-prediction literatures usually conflate ---
\emph{modelling} the conditional mean (judged by likelihood) and
\emph{calibrating} the predictive tail (optionally judged by a downstream score
such as the continuous ranked probability score, CRPS). We call the resulting
recipe \textbf{model first, conform last}, and show that it matches a dedicated
CRPS specialist on CRPS while \emph{improving} held-out log-likelihood. In a fair
rolling one-step-ahead comparison on 500 \proglang{FRED} change-series against
\emph{eight} distributional baselines --- \pkg{AutoARIMA}, \pkg{AutoETS},
\pkg{statsmodels} \pkg{SARIMAX} and exponential smoothing, a \pkg{GARCH}$(1,1)$
model with Student-$t$ innovations, a \pkg{NeuralForecast} model with a
Student-$t$ head, and \pkg{AutoARIMA} paired with split-conformal and
adaptive-conformal residuals --- the general forecaster \textbf{wins the
per-series log-likelihood race against all eight}. The decisive test is the
\pkg{GARCH}-$t$ model, the classical heavy-tail / volatility-clustering state of
the art: our forecaster still edges it on the continuous subset (mean continuous
log-likelihood $2.86$ versus $2.72$), confirming the heavy-tail advantage
\emph{against the heavy-tail specialist}. On CRPS the picture is mixed, as
expected: we beat the smoothing, \pkg{GARCH} and neural baselines, roughly tie the
\pkg{ARIMA} family, and lose to the CRPS-optimised conformal variants. We further
contribute a \emph{mean-preserving lattice projection} for series that revisit
exact values, an online \emph{coordinate} (Yeo--Johnson) search, and a committed
\emph{martingale} model whose volatility clock is a time-changed Brownian motion.
The whole library is implemented twice --- pure \proglang{Python} and
zero-dependency \proglang{JavaScript} --- and verified identical to $10^{-6}$ on
roughly 90{,}000 probe values, so the same models run server-side or in the
browser via \pkg{Pyodide}.
}

\begin{document}

\maketitle

\begin{abstract}
\skatersabstract
\end{abstract}

\noindent\textbf{Keywords:} probabilistic forecasting, time series, proper
scoring rules, conformal prediction, online learning, \proglang{Python},
\proglang{JavaScript}, \pkg{skaters}.

%% -------------------------------------------------------------------------
\section{Introduction} \label{sec:intro}

Most deployed forecasters return a point and, perhaps, an interval bolted on
afterwards. Yet the quantity a decision-maker actually needs is the
\emph{predictive distribution}: the mean for an estimate, the spread for risk,
the tails for position sizing, and a density for likelihood-based evaluation and
combination. Two strands of recent work take distributions seriously but each
gives something up.

Classical state-space and exponential-smoothing models
\citep{hyndman2002state,hyndman2008automatic} are distributional but rarely
\emph{online and composable}: adding a transform usually means re-deriving the
filter. Conditional-heteroscedasticity models \citep{engle1982autoregressive,
bollerslev1986generalized} capture the heavy tails and volatility clustering of
financial data, but commit to a single parametric form and are refit in batch.
At the other extreme, \textbf{conformal prediction}
\citep{vovk2005algorithmic,vovk2019conformal} is distribution-free and online,
but a conformal predictive system outputs a \emph{cumulative distribution
function}, not a density. It is calibrated for coverage and CRPS but
\textbf{cannot be scored on log-likelihood at all}, and an empirical conformal CDF
assigns zero density --- and hence $-\infty$ log-likelihood --- to any value
outside the observed residual range. Modern neural and foundation models
\citep{salinas2020deepar,ansari2024chronos,das2024timesfm,woo2024moirai} are
powerful distributional learners, but are built for global cross-series training
on hardware accelerators, not for a single univariate stream running in a
browser.

\pkg{skaters} takes a third route. Every node --- transform, ensemble, or leaf
--- is a function \code{(y, state) -> ([Dist], state)} that consumes one
observation and emits a list of \code{Dist} objects, one per forecast horizon,
where a \code{Dist} is a weighted Gaussian mixture. Because the type is closed
under composition, an entire model is a chain of invertible transforms with one
distributional leaf at the bottom, and ensembles are simply skaters that combine
other skaters. This single abstraction yields three contributions developed in
turn:

\begin{enumerate}
\item \textbf{A pluggable objective} (Sections~\ref{sec:trunk}--\ref{sec:leaf}).
  The leaf fits its mixture weights by \emph{a} proper scoring rule. Point it at
  log-likelihood and it models honestly; point its \emph{terminal} copy at CRPS
  while keeping a likelihood-weighted trunk, and you get \textbf{model first,
  conform last} --- competitive with a CRPS specialist on CRPS \emph{and} better
  on likelihood.
\item \textbf{Distributional structure that survives composition}
  (Sections~\ref{sec:lattice}--\ref{sec:doob}). A \emph{mean-preserving lattice
  projection} places atoms on the exact values a series revisits without moving
  the ensemble mean and while vanishing on continuous data; an online
  \emph{coordinate} search learns whether a series is simple in a log, root, or
  linear coordinate; and a committed \emph{martingale} model learns only a
  volatility clock.
\item \textbf{An honest benchmark against eight distributional baselines}
  (Section~\ref{sec:experiments}), including the heavy-tail (\pkg{GARCH}-$t$) and
  neural (\pkg{NeuralForecast}) state of the art, scored on held-out
  log-likelihood and CRPS in a single, identical \code{Dist}-based protocol.
\end{enumerate}

The unifying premise is the No-Free-Lunch theorem \citep{wolpert1997no} taken
seriously: every ``method'' is a prior, and averaged over all series none
dominates. The right response is not to pick one but to make the priors
explicit, weight them online, and bound the cost of a wrong prior by how fast the
weighting abandons it.

%% -------------------------------------------------------------------------
\section{The \code{Dist} abstraction and composable transforms} \label{sec:dist}

A \code{Dist} is a weighted mixture $\sum_i w_i\,\mathcal N(\mu_i,\sigma_i^2)$
with $\sigma_i>0$ and $\sum_i w_i = 1$. It exposes \code{mean}, \code{std},
\code{pdf}, \code{cdf}, \code{logpdf}, \code{quantile}, and a closed-form
\code{crps}, together with the affine maps (\code{shift}, \code{scale},
\code{affine}) and a moment-matching \code{prune} needed to propagate it through
transform inverses and to bound component growth in ensembles. The Gaussian
mixture is a deliberate choice: it is closed under affine maps and mixing, admits
a closed-form CRPS through the elementary identity
$\mathbb{E}\,|\mathcal N(m,s^2)| = 2 s\,\phi(m/s) + m\,(2\Phi(m/s)-1)$, and
approximates heavy tails arbitrarily well as a scale mixture.

A \emph{transform} is an online bijection given by a \code{forward} (scalar in,
scalar out) and an \code{inverse\_k} that maps $k$ \code{Dist}s in the
transformed space back to the original. Differencing, exponential smoothing,
drift, Holt's linear trend, fractional differencing, autoregression (via online
recursive least squares), \pkg{GARCH}-style scaling, seasonal differencing,
power and Yeo--Johnson coordinates \citep{box1964analysis,yeo2000new}, and
standardisation are all transforms. \emph{Conjugation} composes a transform with
an inner skater,
\[
y \;\xrightarrow{\;T\;}\; y' \;\to\; \text{inner} \;\to\; D'
  \;\xrightarrow{\;T^{-1}\;}\; D ,
\]
and because every step is online and the inverse propagates the full mixture,
multi-step uncertainty accumulates correctly --- for example variance grows as
$\sum_{h}\sigma_h^2$ under differencing.

\paragraph{Numerical note.}
Mixture variance is computed in centred form,
$\operatorname{Var} = \sum_i w_i\,(\sigma_i^2 + (\mu_i-\bar\mu)^2)$, rather than
$\mathbb E[X^2]-\mathbb E[X]^2$. The latter catastrophically cancels when a tight
component sits at a large mean, silently producing $\sigma=0$ and an invalid
predictive. Running moments use Welford's recurrence \citep{welford1962note}.

%% -------------------------------------------------------------------------
\section{The terminal-leaf ensemble: model first} \label{sec:trunk}

Bayesian model averaging over a heterogeneous pool
\citep{hoeting1999bayesian,raftery2005using} combines full predictive
distributions, which preserves the combined mean and variance but \emph{washes
out higher moments}: a heavy-tailed leaf used inside the pool collapses toward a
Gaussian shape at the output, because a mixture of differently located Gaussians
is dominated by the spread of their means. We therefore use the candidate models
only as \textbf{mean forecasters}, weight them online by predictive likelihood
(with a learning-rate shrinkage and a per-depth complexity penalty, in the
spirit of gradient boosting), combine their means, and model the distribution of
the \emph{combined residual} with a \textbf{single terminal leaf}:
\[
 y \;\to\; \hat\mu = \sum_i w_i\,\hat\mu_i, \qquad
 r = y-\hat\mu \;\to\; \text{leaf} \;\to\; D, \qquad
 \hat D = D \text{ shifted by } \hat\mu .
\]
Because exactly one leaf reaches the output, its shape --- heavy tails, learned
online --- survives undiluted. This is the \emph{model first} half of the recipe:
a likelihood-weighted trunk whose only job is to get the conditional mean right.
The candidate pool spans exponential smoothing, differencing, drift, Holt's
linear trend, autoregression, grouped long-lag autoregression, fractional
differencing, a seasonal block, a Yeo--Johnson coordinate grid
(Section~\ref{sec:coord}), and a fast/slow two-systems block.

%% -------------------------------------------------------------------------
\section{Metric-agnostic leaves: conform last} \label{sec:leaf}

The leaf is itself a fixed Gaussian \emph{scale mixture} --- components
$\mathcal N(0, c_k\,\hat\sigma)$ over a fixed log-spaced dictionary $\{c_k\}$,
with $\hat\sigma$ an exponentially weighted scale --- whose only learned
quantities are the mixture weights $\{w_k\}$ on the simplex. A Gaussian scale
mixture is a Student-$t$ in the limit, so the leaf approximates heavy tails by
construction while remaining a plain \code{Dist}. The weights are fit online by
\textbf{a proper scoring rule} \citep{gneiting2007strictly}:
\begin{itemize}
\item \code{likelihood} --- recency-weighted expectation--maximisation on the
  component responsibilities;
\item \code{crps} --- exponentiated-gradient descent on the simplex minimising the
  closed-form mixture CRPS \citep{matheson1976scoring}; the gradient is exact,
  since the mixture CRPS is a finite sum of $\mathbb E\,|\mathcal N(m,s^2)|$ terms.
\end{itemize}

Placing the CRPS leaf at the \emph{terminal} of a likelihood-weighted trunk gives
\textbf{model first, conform last}. The trunk's job is to model, and the honest
objective for modelling is likelihood --- a proper, tail-sensitive score that is,
up to a constant, the Kelly log-growth criterion \citep{kelly1956new}; corrupting
it to chase CRPS would repeat the conformal mistake of discarding the density.
The tail's job is to be shaped, and that is the only place a downstream metric
(CRPS) should get a vote. Section~\ref{sec:ablations} shows this configuration
dominates a pure-likelihood policy on CRPS \emph{and} on likelihood: a CRPS-fit
leaf is more outlier-robust, so it even generalises slightly better in
log-likelihood. The objective is a knob, not a fixed cost.

%% -------------------------------------------------------------------------
\section{The lattice projection} \label{sec:lattice}

Administrative and grid-quoted series --- policy rates, posted prices, pegged
exchange rates --- revisit a small set of exact values, and a continuous
predictive density cannot place mass there. We add a \emph{projection}: a
recency-weighted frequency table of exact values, with near-Dirac atoms on the
values revisited above a noise floor, each carrying its observed frequency as
probability. The atom is \textbf{mean-preserving}: the continuous part is
recentred so the atoms add \emph{mass} without moving the ensemble mean.

Two properties make this a safe default. First, on continuous data nothing is
revisited, no atom fires, and the projection vanishes identically --- it is free
when unused. Second, unlike a last-value spike it captures values that recur
often but never consecutively (a rate that repeatedly returns to a round number).
The pull of the level (its tendency to persist) is left to the trunk, where the
random-walk candidate earns its weight by likelihood, so the projection remains a
pure statement about mass, not mean.

%% -------------------------------------------------------------------------
\section{Coordinate learning} \label{sec:coord}

Whether a series is simple in a linear, log/multiplicative, or root coordinate is
itself learnable. We add a coarse grid of \textbf{Yeo--Johnson} $\lambda$
candidates --- the signed Box--Cox family \citep{box1964analysis,yeo2000new},
defined on all of $\mathbb R$ --- to the pool, so the ensemble learns the
coordinate online rather than committing up front. This is No-Free-Lunch-safe: a
wrong coordinate is simply down-weighted by likelihood. The same mechanism
expresses a non-negativity prior at $\lambda=0$ (the log coordinate), though we
note that because the \code{Dist} is a Gaussian mixture the inverse uses the
delta-method linearisation, so non-negativity is approximate rather than exact.

%% -------------------------------------------------------------------------
\section{A martingale with a learned volatility clock} \label{sec:doob}

For near-martingale \emph{levels} we provide a committed model, \code{doob}: the
mean is pinned to the last value (a driftless martingale) and only the volatility
clock is learned, as a Bayesian average over martingale predictives that differ
in their volatility model (constant, \pkg{GARCH}, slowly varying, heavy-tailed).
Because every candidate shares the same mean, plain averaging blends the clocks
\emph{without} washing out kurtosis --- the committed mean is exactly what makes
model averaging the right tool here, in contrast to Section~\ref{sec:trunk}. By
the Dambis--Dubins--Schwarz theorem
\citep{dambis1965decomposition,dubins1965continuous,doob1953stochastic} a
continuous martingale is a time-changed Brownian motion, so the model is
literally ``Brownian motion on a stochastic clock''. It beats the general
forecaster on near-martingale levels by committing the mean and spending its
capacity on the clock; on genuinely mean-reverting series the martingale prior is
wrong and it gives ground --- a deliberately sharp instrument, fed \emph{levels}
rather than pre-differenced changes.

%% -------------------------------------------------------------------------
\section{Experiments} \label{sec:experiments}

\paragraph{Universe.}
To avoid a hand-picked series list we draw from the Federal Reserve Economic Data
\citep{fred}, taking cached daily series and auto-transforming each to its
one-step \emph{changes} (log-difference for strictly positive levels, else first
difference) --- the stationary-ish innovation stream where heavy tails and regime
shifts live, and which keeps log-likelihood comparable across series of different
scale. We keep the last $1000$ changes per series and score the last $300$ as a
rolling one-step-ahead test window after a burn-in.

\subsection{Against eight distributional baselines} \label{sec:sota}

On $500$ such series we compare the general forecaster \code{laplace} against
every distributional one-step baseline we could assemble:
\pkg{statsforecast}'s \citep{garza2022statsforecast} \pkg{AutoARIMA} and
\pkg{AutoETS}; \pkg{statsmodels} \pkg{SARIMAX}$(1,0,1)$ and simple exponential
smoothing, which emit \emph{exact} closed-form Gaussian predictives;
a constant-mean \pkg{GARCH}$(1,1)$ model with Student-$t$ innovations
\citep{bollerslev1986generalized,bollerslev1987conditionally} from the \pkg{arch}
package, the classical heavy-tail / volatility-clustering state of the art; a
\pkg{NeuralForecast} \citep{olivares2022neuralforecast} multilayer perceptron
with a Student-$t$ distribution head; and \pkg{AutoARIMA} paired with conformal
residual quantiles \citep{vovk2019conformal} in both a split-conformal and an
adaptive-conformal (ACI, \citealp{gibbs2021adaptive}) variant.

The protocol is a \emph{fair rolling one-step-ahead}: each baseline is fit on an
expanding or rolling window with periodic refit, and every method --- ours and
theirs --- is turned into the \emph{same} \code{Dist} and scored on held-out
log-likelihood \emph{and} CRPS over the same test window. The parametric
Student-$t$ predictives (\pkg{GARCH}, \pkg{NeuralForecast}) are represented as
finite Gaussian scale mixtures, matching the analytic $t$ density to $\sim
10^{-3}$. We report \textbf{per-series} win-rates (the fraction of series on
which \code{laplace} scores better), and additionally restrict to the $309$
\emph{continuous} series --- those with under $5\%$ exactly-repeating one-step
changes --- since on repeating/grid series the lattice projection
(Section~\ref{sec:lattice}) gives a large but metric-specific advantage that
would otherwise dominate the aggregate. Table~\ref{tab:sota} reports the result.

\begin{table}[t!]
\centering
\begin{tabular}{lccc}
\toprule
Baseline & LL (all / cont.) & CRPS (all / cont.) & Mean LL (cont.) \\
\midrule
\pkg{AutoARIMA}             & \textbf{78 / 64\,\%} & 51 / 47\,\% & 2.09 \\
\pkg{AutoETS}               & \textbf{92 / 88\,\%} & 68 / 67\,\% & 2.20 \\
\pkg{SARIMAX} (statsmodels) & \textbf{82 / 72\,\%} & 53 / 49\,\% & 2.13 \\
ETS (statsmodels)           & \textbf{93 / 89\,\%} & 69 / 68\,\% & 2.12 \\
\pkg{GARCH}$(1,1)$-$t$      & \textbf{79 / 67\,\%} & 76 / 66\,\% & \textbf{2.72} \\
NF-Student$t$               & \textbf{100 / 100\,\%} & 78 / 76\,\% & 0.82 \\
\pkg{AutoARIMA} + conformal & \textbf{84 / 75\,\%} & 37 / 29\,\% & 1.47 \\
\pkg{AutoARIMA} + ACI       & \textbf{87 / 80\,\%} & 36 / 28\,\% & 1.89 \\
\bottomrule
\end{tabular}
\caption{Per-series win-rate of \code{laplace} versus each baseline on 500
\proglang{FRED} change-series, all and on the 309 continuous series. LL counts
series where \code{laplace} has the higher held-out log-likelihood; CRPS where it
has the lower CRPS. The final column is each baseline's mean continuous
log-likelihood (\code{laplace}: \textbf{2.855}). Higher is better for LL; lower for
CRPS.}
\label{tab:sota}
\end{table}

\paragraph{The likelihood race.}
\code{laplace} wins the per-series log-likelihood race against \emph{all eight}
baselines, on the full universe and on the continuous subset, including the two
exact-Gaussian references (\pkg{SARIMAX}, statsmodels ETS) for which there is no
band-reconstruction approximation to blame.

\paragraph{\pkg{GARCH}-$t$ is the decisive test.}
The \pkg{GARCH}$(1,1)$-$t$ model is purpose-built for the heavy tails and
volatility clustering of financial change-series, so it is \emph{supposed} to be
the hardest opponent on log-likelihood --- and it is. It is the tightest race
($67\%$ of continuous series), and its mean continuous log-likelihood, $2.72$, is
closest to ours, $2.855$ --- a genuine but \emph{narrow} edge of about $0.14$
nats per observation. That the advantage survives \emph{against the heavy-tail
specialist} is the result that matters: the general forecaster captures the same
conditional heteroscedasticity without committing to a parametric volatility law.

\paragraph{The neural result is not a scalp.}
\code{laplace} beats NF-Student$t$ on $100\%$ of series, but we are deliberately
loud about what this does and does not mean. NF here is \pkg{NeuralForecast} run
in the \emph{matched online univariate one-step} protocol --- a tiny per-series
multilayer perceptron with periodic refit --- which is its \textbf{weak} regime.
Neural and foundation forecasters \citep{salinas2020deepar,
olivares2022neuralforecast} are built for \emph{global cross-series} training,
where they amortise structure across thousands of related series. Our result says
nothing about \pkg{DeepAR}-style models in that home setting, and we do not claim
it does; it says only that a single short univariate stream does not give a neural
density head enough to beat an online ensemble. (We verified the NF densities are
well formed --- median log-likelihood $0.50$, none floored --- so the $100\%$ is a
real calibration gap, not a degenerate score.)

\paragraph{CRPS.}
As ever, CRPS tells a mixed story. \code{laplace} beats the smoothing baselines
(\pkg{AutoETS}, ETS), \pkg{GARCH}-$t$ and NF on CRPS, roughly ties the
\pkg{ARIMA} family (\pkg{AutoARIMA} $51\%$, \pkg{SARIMAX} $53\%$), and
\emph{loses} to the conformal variants ($37\%$, $36\%$), which are CRPS-optimised
by construction. This is exactly the thesis of Sections~\ref{sec:leaf}: likelihood
is the metric where a faithful density wins; CRPS is conformal's home turf, where
we are competitive but not dominant. A practitioner who truly only cares about
CRPS can repoint the leaf (Section~\ref{sec:ablations}); a practitioner who cares
about the density --- for risk, sizing, or combination --- should prefer the
honest log-likelihood column.

\paragraph{Scope.}
Five hundred change-series, one-step horizon. Prophet, multi-step horizons and raw
levels are left for later; zero-shot foundation models get their own protocol in
Section~\ref{sec:foundation}.

\subsection{Against zero-shot foundation models} \label{sec:foundation}

Pretrained time-series foundation models use a fundamentally different protocol:
they are applied \emph{zero-shot}, conditioning on a context window with no
fitting. We evaluate four ---
Chronos-Bolt \citep{ansari2024chronos}, TimesFM~2.5 \citep{das2024timesfm},
Moirai \citep{woo2024moirai}, and Lag-Llama \citep{rasul2023lagllama} --- on 120
change-series (69 continuous). At each step the model receives a fixed
$256$-length window of the preceding changes and predicts the next change, with
all windows for a series batched into one inference call. Each predictive is
turned into the same \code{Dist} --- the native Student-$t$ head for
Moirai/Lag-Llama, a sample/quantile reconstruction for Chronos-Bolt/TimesFM ---
and \code{laplace} is re-scored on the identical window. Table~\ref{tab:foundation}
reports per-series win-rates.

\begin{table}[t!]
\centering
\begin{tabular}{llccc}
\toprule
Model & density & LL (all / cont.) & CRPS (all / cont.) & Mean LL (cont.) \\
\midrule
Moirai (1.1-R-small)  & native $t$    & 98 / \textbf{97}\,\% & 94 / 93\,\% & 0.92 \\
Lag-Llama             & native $t$    & 67 / \textbf{99}\,\% & 65 / 94\,\% & 0.39 \\
Chronos-Bolt (small)  & quantile$^{*}$ & 76 / \textbf{100}\,\% & 67 / 88\,\% & $-1.96$ \\
TimesFM (2.5-200M)    & quantile$^{*}$ & 75 / \textbf{100}\,\% & 58 / 72\,\% & $-1.18$ \\
\bottomrule
\end{tabular}
\caption{Per-series win-rate of \code{laplace} versus four zero-shot foundation
models on 120 change-series (69 continuous). $^{*}$Chronos-Bolt/TimesFM
log-likelihood is reconstructed from a quantile head and is tail-limited; read
their CRPS as the fairer signal. \code{laplace} mean continuous logpdf:
\textbf{1.51}.}
\label{tab:foundation}
\end{table}

\emph{On the continuous series, \code{laplace} beats all four foundation models
zero-shot on both metrics} (97--100\% LL, 72--94\% CRPS). The native-density
models (Moirai, Lag-Llama) are the meaningful likelihood opponents and the
closest, but still lose per-series. On repeat-heavy/grid series the foundation
models are competitive or better --- Lag-Llama even beats \code{laplace} on mean
log-likelihood over the \emph{full} universe ($6.54$ vs $4.52$) by placing tight
mass on revisited values, the same effect as the lattice projection --- which is
why the continuous split matters. The honest scope is narrow: zero-shot, no refit,
fixed context, and forecasting the change stream rather than the levels these
models were chiefly trained on.\footnote{Per-series \emph{fine-tuning} does not
rescue the foundation models on this task: naively fine-tuning Lag-Llama on a
single series' history (5 epochs) \emph{catastrophically overfits}, collapsing its
held-out log-likelihood from $+1.5$ to below $-100$ while taking $\sim$15$\times$
longer. This is the expected failure mode --- adapting a pretrained model to one
short univariate stream is not its design regime; zero-shot (or domain-level
fine-tuning across many series) is the intended use. We therefore report the
zero-shot result as the headline.}

%% -------------------------------------------------------------------------
\section{Ablations: repointing the objective} \label{sec:ablations}

A companion sweep on $2{,}501$ \proglang{FRED} series isolates the effect of the
terminal-leaf objective against the \pkg{crepes} conformal predictive system
(given its three best calibration windows and scored on its own CDF via the
pinball decomposition of CRPS). Because the conformal opponent there uses only a
naive mean, we read this as an ablation of our own objective knob rather than a
state-of-the-art claim. Reporting per-policy and de-correlating to one vote per
correlated family (195 families):

\begin{table}[h!]
\centering
\begin{tabular}{lccc}
\toprule
Forecaster (identical structure) & CRPS, raw & CRPS, family & Mean LL \\
\midrule
likelihood trunk + likelihood leaf              & 82.6\,\% & 48.7\,\% & 2.96 \\
likelihood trunk + CRPS leaf (\emph{model first, conform last}) & 89.5\,\% & \textbf{60.1\,\%} & \textbf{3.01} \\
bare CRPS leaf (no trunk)                        & 80.8\,\% & 60.3\,\% & 2.96 \\
\bottomrule
\end{tabular}
\caption{Effect of the terminal-leaf objective. A general likelihood policy is
essentially even with conformal on CRPS over independent families ($48.7\%$);
repointing only the terminal leaf to CRPS lifts that to $60.1\%$ \emph{and} raises
log-likelihood, matching the dedicated CRPS specialist while modelling honestly.}
\label{tab:objective}
\end{table}

Three further observations motivate the radically small public API.
\textbf{Priors wash out:} across $25$ series with $\geq 2000$ changes the
conventional prior-flavoured policies (trend, long-memory, drift, minimax,
fast/slow) sit within $0.0006$ nats of the uniform-prior general forecaster on
the second half of each series, and the first-half winner beats the general
forecaster on the second half only $44\%$ of the time --- worse than chance.
\textbf{Conform-last is near-free:} on idealised light-tailed processes the CRPS
leaf costs $0.005$--$0.007$ nats versus the likelihood leaf, while on heavy-tailed
and stochastic-volatility processes it \emph{helps} (up to $+0.04$); real data
lives in the second regime. \textbf{The lattice projection is free when unused.}
Together these justify exposing exactly \emph{one} general forecaster
(uniform prior, CRPS leaf, lattice projection and coordinate learning on by
default) plus \emph{one} committed martingale specialist; the remaining
structural ideas do not pay their way out of sample.

%% -------------------------------------------------------------------------
\section{Discussion} \label{sec:discussion}

The unifying theme is No-Free-Lunch \citep{wolpert1997no} taken seriously. Every
``method'' is a prior; averaged over all series none dominates. The correct
response is not to choose one but to make the priors explicit, weight them online,
and bound the cost of a wrong prior by how fast the weighting abandons it. Read
that way, the named methods of classical forecasting are convenience bundles, and
the honest object is a single adaptive ensemble with two orthogonal terminal knobs
--- \emph{which proper scoring rule the leaf optimises}, and \emph{whether a
lattice projection is active} --- neither of which is a prior over the trunk. The
empirical headline supports the design: against the strongest distributional
opponents available on CPU, including the heavy-tail and neural state of the art,
the general forecaster wins the likelihood race, and where it does not lead ---
CRPS --- the deficit is confined to methods explicitly optimised for that metric
and is recoverable by repointing a single knob.

%% -------------------------------------------------------------------------
\section{Heritage and reproducibility} \label{sec:heritage}

\pkg{skaters} is a from-scratch rewrite distilling ideas from \pkg{timemachines}
\citep{cotton2021timemachines}, an earlier competition-tested online-forecasting
package, and builds on years of running live distributional-prediction contests
at Microprediction, where forecasts are scored continuously as full
distributions. The library is implemented twice --- pure \proglang{Python} and
zero-dependency \proglang{JavaScript} --- and a parity suite checks roughly
$90{,}000$ probe values to $10^{-6}$ on every run, so results are reproducible
across both runtimes and in the browser via \pkg{Pyodide} \citep{pyodide2021}.
The benchmark of Section~\ref{sec:experiments} is a single parallel script
(\code{benchmarks/sota\_study.py}) that scores every method through the identical
\code{Dist} interface; given a \proglang{FRED} key and the optional baselines
(\pkg{statsforecast}, \pkg{arch}, \pkg{statsmodels}, \pkg{neuralforecast}) it
reproduces Table~\ref{tab:sota} end to end.

%% -------------------------------------------------------------------------
\bibliographystyle{plainnat}
\bibliography{skaters}

\end{document}
